\chapter{Gene duplication}
\label{ch:duplication}

\section{Networks that copy themselves}

The two rules of the preceding chapters each had a newcomer reach into the
existing network and \emph{choose} where to attach --- by surfing, or by closing
a triangle. Biology offers a third rule, and in some ways a more primitive one:
the network grows by \emph{copying} a node that is already there. The newcomer
does not pick its neighbors at all. It inherits them, wholesale, from a parent.

The setting is the inner machinery of the living cell. Genomes evolve in good
part through \emph{duplication} --- a gene, or sometimes an entire genome, is
accidentally copied --- followed by \emph{divergence}, the slow accumulation of
mutations that make the two copies drift apart. That duplication is a major
engine of evolution was argued classically by \citet{ohno1970}. To turn this into
a network statement we use the \emph{protein interaction network}: each node is a
protein (the product of a gene), and a link joins two proteins that physically
interact --- that bind to each other to do their job. Now watch what a gene
duplication does to this network. Right after the gene is copied, the cell makes
two identical proteins, so the new protein interacts with \emph{exactly the same
partners} as the old one: the duplicate node arrives already wired to all the
neighbors of its ancestor (Figure~\ref{fig:dup-diverge}). Mutations then add a few
interactions and, more often,
remove some, letting the twins diverge \citep{wagner2001,sole2002}. Both steps
are thoroughly local --- the copy needs to know only its parent and its parent's
neighbors, never the rest of the genome.

That moment just after the copy, when two proteins have \emph{exactly} the same
partners, is worth marking now because it turns out to matter far beyond the degree
statistics this chapter is about. Two nodes with identical neighborhoods are an exact
symmetry of the network, and Chapter~\ref{ch:quantum} shows that such a pair forces a
quantum walker to stay put: it pins an eigenvector of the walk onto the twins and
nowhere else. Duplication, in other words, does not merely shape how many links a
protein has. It plants a spectral signature that survives in the measured
interactome, and the same chapter finds real protein networks to be the most strongly
localized in this book.

\begin{figure}[t]
\centering
\begin{adjustbox}{max width=\linewidth}\begin{tikzpicture}[
  nd/.style={circle,draw=black!65,fill=black!4,minimum size=5.5mm,inner sep=0pt,font=\scriptsize},
  cp/.style={circle,draw=accent,fill=black!30,minimum size=5.5mm,inner sep=0pt,font=\scriptsize},
  lnk/.style={draw=black!45,line width=0.8pt},
  newl/.style={draw=accent,line width=1.3pt},
  gone/.style={draw=black!22,line width=0.7pt,dotted}]
  % panel (a) duplication
  \begin{scope}
    \node[nd] (i)  at (-0.2,0.6)  {$i$};
    \node[cp] (ip) at (-0.2,-0.6) {$i'$};
    \node[nd] (a)  at (1.5,1.05)  {};
    \node[nd] (b)  at (1.75,0.0)  {};
    \node[nd] (c)  at (1.5,-1.05) {};
    \draw[lnk]  (i)--(a); \draw[lnk] (i)--(b); \draw[lnk] (i)--(c);
    \draw[newl] (ip)--(a); \draw[newl] (ip)--(b); \draw[newl] (ip)--(c);
    \draw[newl] (i)--(ip);
    \node[font=\footnotesize] at (0.75,-1.85) {(a) duplication};
  \end{scope}
  \node at (3.25,0.0) {\large $\Rightarrow$};
  % panel (b) divergence
  \begin{scope}[xshift=5.7cm]
    \node[nd] (i2)  at (-0.2,0.6)  {$i$};
    \node[cp] (ip2) at (-0.2,-0.6) {$i'$};
    \node[nd] (a2)  at (1.5,1.05)  {};
    \node[nd] (b2)  at (1.75,0.0)  {};
    \node[nd] (c2)  at (1.5,-1.05) {};
    \draw[lnk]  (i2)--(a2);
    \draw[newl] (ip2)--(b2);
    \draw[lnk]  (i2)--(c2);
    \draw[newl] (i2)--(ip2);
    \draw[gone] (ip2)--(a2);
    \draw[gone] (i2)--(b2);
    \draw[gone] (ip2)--(c2);
    \node[font=\footnotesize] at (0.75,-1.85) {(b) divergence};
  \end{scope}
\end{tikzpicture}\end{adjustbox}
\caption{\textbf{Duplication and divergence.} (a) A gene duplication adds a copy
$i'$ (gray) wired to every interaction partner of the ancestor $i$, plus a link
between the twins when the protein is self-interacting. (b) Divergence then prunes
the now-redundant interactions (dotted), each partner kept by only one twin, so
the copies drift apart.}
\label{fig:dup-diverge}
\end{figure}

The same picture reaches beyond biology. A new web page is often made by copying
another page's hyperlinks and then editing the list --- duplication as copying,
divergence as editing. Following \citet{vazquez2003} and \citet{vazquez2003pin},
this chapter shows that copy-and-diverge, like the rules before it, produces an
effective ``rich get richer,'' a power-law degree distribution, the clustering
hierarchy, and --- as actually measured in protein networks \citep{maslov2002} ---
\emph{disassortative} correlations. It also reveals something new: a
\emph{multifractal} degree distribution, which we will unpack when we reach it, a
fingerprint left by the act of inheritance.

\section{Duplication and divergence, in their simplest form}

Start with the simplest version, treating duplication and divergence as
independent random events. The single idea that drives everything is this: a node
gains a new link whenever one of its \emph{neighbors} is duplicated, because the
neighbor's copy inherits a link back to it. A well-connected node has many
neighbors, so it has many chances to be on the receiving end of a duplication ---
and there, already, is the rich getting richer. The box makes this precise and
reads off the degree distribution; the conclusion follows it.

\begin{calculation}{degree growth and the power law from duplication}
In a mean-field (average-quantities) description, with increasing the number of
nodes $n$ the degree of node $i$ evolves as
\begin{equation}
  \frac{\partial k_i}{\partial n}
    = \nu_D\, k_i + \nu_C\,(n - k_i) - \nu_L\, k_i,
  \label{eq:dup-master}
\end{equation}
with $\nu_D$, $\nu_C$, $\nu_L$ the per-node rates of duplication, random link
creation, and link loss. The first term is the key one: it is proportional to
$k_i$ because each of $i$'s $k_i$ neighbors, if duplicated, hands $i$ a new link.
Each duplication adds one node, so $\nu_D = 1/n$; writing $\nu_C=\mu_0/n$ and
$\nu_L=\mu_1/n$ reduces the equation to
\begin{equation}
  n\,\frac{\partial k_i}{\partial n} = \mu_0 + (1-\mu_1)\, k_i,
  \label{eq:dup-reduced}
\end{equation}
whose growth rate is linear in the degree. Integrating it gives the degree of
node $i$ as the network grows,
\begin{equation}
  k_i(n) = \left(k_i(n_i) + \frac{\mu_0}{1-\mu_1}\right)
           \left(\frac{n}{n_i}\right)^{\beta} - \frac{\mu_0}{1-\mu_1},
  \label{eq:dup-growth}
\end{equation}
where $n_i$ and $k_i(n_i)$ are the network size and the node's degree at the
moment it was added, and the exponent is
\begin{equation}
  \beta = 1 - \mu_1 \qquad (\mu_1 < 1),
  \label{eq:dup-beta}
\end{equation}
the condition $\mu_1<1$ keeping the network from emptying out. To turn this growth
law into a degree distribution, note that node $i$ has degree above $k$ exactly
when its age and birth degree make Eq.~\eqref{eq:dup-growth} exceed $k$,
\begin{equation}
  P(k_i > k) = \Theta\!\left[\left(k_i(n_i) + \frac{\mu_0}{1-\mu_1}\right)
           \left(\frac{n}{n_i}\right)^{\beta} - \frac{\mu_0}{1-\mu_1} > k\right],
  \label{eq:dup-cumulative}
\end{equation}
with $\Theta[\cdot]$ the step function, $1$ when its bracket holds and $0$
otherwise.
Both its birth size $n_i$ and its birth degree $k_i(n_i)$ are random: for nodes
added at a constant rate $n_i$ is uniform, $P(n_i{=}n)=1/n$, while a node inherits
its ancestor's degree, so $P[k_i(n_i){=}k']=p_{k'}$ in the stationary state.
Averaging over both,
\begin{equation}
\adjustbox{max width=\linewidth}{$\displaystyle
  P(k_i > k) = \sum_{k'} p_{k'} \int_1^n \frac{dn_i}{n}\,
    \Theta\!\left[\left(k' + \frac{\mu_0}{1-\mu_1}\right)
           \left(\frac{n}{n_i}\right)^{\beta} - \frac{\mu_0}{1-\mu_1} > k\right]
$}
  \label{eq:dup-average}
\end{equation}
Differentiating this cumulative distribution and taking $n\gg1$ leaves,
for nodes added at a constant rate, a power-law degree distribution
\begin{equation}
  p_k = \frac{\partial P(k_i > k)}{\partial k}
      \sim \left(\frac{\mu_0}{1-\mu_1} + k\right)^{-\gamma}, \qquad
  \gamma = 1 + \frac{1}{1-\mu_1}.
  \label{eq:dup-pk}
\end{equation}
\end{calculation}

The result is once again a power law, Eq.~\eqref{eq:dup-pk}, and its exponent is
controlled entirely by $\mu_1$, the rate at which divergence removes links: the
more aggressively the duplicates shed interactions, the steeper the tail and the
fewer the giant hubs. As in the previous two chapters, the power law can be
traced to a single term that grows in proportion to degree --- here the
duplication term of Eq.~\eqref{eq:dup-reduced} --- and without it the distribution
would decay as a featureless exponential. Copying, all by itself, is a
rich-get-richer engine.

\section{Cliquishness from self-interaction}

If divergence added and removed links purely at random, this model would make no
triangles, and its cliquishness would fade to nothing in a large network. But
biology supplies one specific local source of triangles. Some proteins are
\emph{self-interacting} --- they bind to copies of themselves. When the gene for
such a protein is duplicated, the ancestor and its new twin can bind to
\emph{each other}, because each still carries the self-binding site
\citep{wagner2001}. That single interaction, laid down at the moment of
duplication, is the seed of a triangle. The short box turns this into the
clustering law.

\begin{calculation}{the clustering hierarchy from self-interaction}
Let a duplicated self-interacting protein bond to its ancestor with probability
$q_v$. Then whenever a neighbor of node $i$ is duplicated, $i$ gains not only the
copy as a new neighbor but, with probability $q_v$, a link between two of its
neighbors (the ancestor and the copy). Hence the number of links $e_i$ between the
neighbors of node $i$ change as
\begin{equation}
\frac{\partial e_i}{\partial n} \approx q_v\,\frac{\partial k_i}{\partial n}.
\end{equation}
Integrating, assuming $e_i=0$ when $k_i=1$, we get $e_i = q_v (k_i - 1)$.
Inserting this result into the clustering definition,
\begin{equation}
  \langle c\rangle_k = \frac{2 e(k)}{k(k-1)} \approx \frac{2 q_v}{k}.
  \label{eq:dup-ck}
\end{equation}
\end{calculation}

So cliquishness again falls as one over the degree, Eq.~\eqref{eq:dup-ck} --- the
clustering hierarchy of the introduction, reached now by a third independent
route. There is one telling refinement: because removing an interaction also
destroys triangles, and the biggest hubs have the most triangles to lose,
divergence makes the clustering of high-degree nodes decay a little \emph{faster}
than one over the degree. That slightly-too-steep falloff is a small signature
that distinguishes duplication from the search and closure mechanisms.

\section{When duplication and divergence are coupled}

In a real protein network the two processes cannot truly be independent. The
network does a job --- it keeps the cell alive --- so losing the wrong interaction
can be lethal, and which links a duplicate keeps is under natural selection. The
classical view of \citet{ohno1970} lets one redundant copy decay freely; the more
modern \emph{subfunctionalization} picture of \citet{force1999} has the two
copies divide the ancestor's duties between them, each keeping part, together
keeping the whole. Either way the fates of the two copies' links are
\emph{coupled}: a shared interaction is far more likely to be retained by one
copy than dropped by both. To capture this, \citet{vazquez2003pin} use the
following rule. Starting from two connected nodes, repeat:
\begin{itemize}
  \item \emph{Duplication.} Pick a node $i$ at random and add a copy $i'$ joined
    to every neighbor of $i$; with probability $q_v$ also add the link $i$--$i'$
    (the self-interaction case).
  \item \emph{Divergence.} For each node $j$ now joined to both $i$ and $i'$,
    delete one of the two redundant links $(i,j)$ or $(i',j)$, at random, with
    probability $1-q_e$.
\end{itemize}
The divergence step touches only the \emph{redundant} pairs just created, and
that is what couples it to duplication. Genomic data showing that duplicates
diverge soon after they are born \citep{wagner2001} justify treating divergence
as finishing before the next duplication, so the node count $n$ doubles as a
clock. A related model with lopsided divergence is studied in \citet{vazquez2003}.

Coupling makes the large-scale behavior considerably richer, in two ways worth
previewing before the box. First, a single number --- the retention parameter
$q_e$ --- decides whether the network stays sparse or grows ever denser, with a
sharp switch between the two regimes. Second, the degree distribution is no
longer a plain power law but a \emph{multifractal} one. A quick word on that
term: a plain power law is described by one exponent, the same at every scale; a
multifractal is described by a whole continuum of exponents, because different
``moments'' of the distribution (the average degree, the average of the squared
degree, and so on) scale with different powers of the network size. The box
records the transition and the moment exponents; the interpretation follows.

\begin{calculation}{the density transition and the multifractal spectrum}
Tracking the mean degree through one duplication event gives
\begin{equation}
\adjustbox{max width=\linewidth}{$\displaystyle
  \langle k\rangle(n+1)
   = \frac{n\langle k\rangle(n) + 2q_v + (2q_e - 1)\langle k\rangle(n)}{n+1}
$}
  \label{eq:dup-meandeg}
\end{equation}
which in the large-network limit switches behavior at $q_e=1/2$: for $q_e<1/2$
the mean degree saturates at the finite value $\langle k\rangle = 2q_v/(1-2q_e)$,
while for $q_e>1/2$ it grows without bound as $N^{2q_e-1}$. A rate-equation
treatment in the manner of \citet{kim2002} gives an attachment rate
$A(k)=(q_v+q_e k)/n$ --- linear in degree again --- and predicts that the $l$-th
moment of the degree distribution scales as $M_l\sim n^{\sigma_l(q_e)}$ with
\begin{equation}
  \sigma_l(q_e) = l\, q_e + 2\left[\left(\frac{1+q_e}{2}\right)^{l} - 1\right].
  \label{eq:dup-moments}
\end{equation}
For each $l$ there is a value $q_l$ at which $\sigma_l$ changes sign, with
$q_1=1/2$ and $q_2 = 2\sqrt{3}-3 \approx 0.46$.
\end{calculation}

The two messages of the box are these. First, Eq.~\eqref{eq:dup-meandeg} shows a
genuine switch at $q_e=1/2$: keep too few links and the network settles into a
sparse steady state; keep enough and it densifies forever. Second, and more
subtle, the moment exponents $\sigma_l$ in Eq.~\eqref{eq:dup-moments} depend on
$l$ in a curved, \emph{nonlinear} way --- the hallmark of a multifractal. Where
the earlier chapters gave clean single power laws, duplication gives a spectrum
of exponents, and the reason is exactly the inheritance at its core: each
newcomer copies the degree environment of its parent, injecting a kind of
heritable heterogeneity that a single exponent cannot capture. The deletion of
links from the copies is what converts that inheritance into multifractality.

\section{Disassortative correlations}

Simulations spanning network sizes from a thousand to a million nodes confirm the
analysis \citep{vazquez2003pin, vazquez2003}: the moment exponents match the multifractal prediction; the clustering
falls with degree as a power law with exponent at or above one (steepening as
link loss grows, just as Eq.~\eqref{eq:dup-ck} and its caveat foretold); and the
average neighbor degree \emph{decreases} with degree, $\langle k_{nn}\rangle(k)
\sim k^{-0.1}$. The duplication network is therefore \emph{disassortative} --- its
hubs are surrounded by poorly-connected partners rather than by each other.

This is the correct sign. Disassortative mixing is precisely what
\citet{maslov2002} measured in the protein interaction network of yeast, and it
files duplication alongside the technological networks of
Chapter~\ref{ch:search} and against the social networks of
Chapter~\ref{ch:friends}. With this third mechanism in hand, the three rules of
the book so far bracket the empirical landscape, as the table summarizes (here
$\beta$ is the rate at which cliquishness falls with degree, and $\alpha$ the
slope of the neighbor-degree curve --- positive for assortative, negative for
disassortative):
\begin{center}
\begin{tabular}{lcc}
\hline
Mechanism & $\langle c\rangle_k \sim k^{-\beta}$ & $\langle k_{nn}\rangle_k \sim k^{\alpha}$ \\
\hline
Connecting neighbors    & $0 < \beta < 1$ & $\alpha > 0$ \\
Random walk             & $\beta = 1$     & $\alpha \le 0$ \\
Duplication--divergence & $\beta \ge 1$   & $\alpha < 0$ \\
\hline
\end{tabular}
\end{center}
Every local rule produces a clustering hierarchy; what differs is the \emph{sign}
of the correlations, which sorts the rules into the assortative (social,
closure-driven) and disassortative (technological and biological, reach-driven)
classes seen in real data.

\section{The unifying principle}

Three different stories --- surfing a network, closing a triangle, copying a node
--- have each delivered the same effective ``rich get richer.'' That is no
accident, and the duplication picture makes the reason simple enough to state in
a single sentence.

Picking a node \emph{at random} plays no favorites: every node is equally likely,
hub or not. But none of the local rules in this book act on a randomly chosen
node. They act on a \emph{neighbor} of one --- the page a surfer steps to, the
friend a friend introduces, the partner a duplicated gene inherits. And the
chance that a particular node $i$ is the neighbor of a randomly chosen node is
\begin{equation}
  \frac{k_i}{\sum_j k_j},
  \label{eq:dup-unifying}
\end{equation}
which is exactly the linear preferential-attachment probability of the
Barab\'asi--Albert model \citep{barabasi1999}. High-degree nodes are
over-represented among the neighbors of a random node for the most elementary of
reasons: they belong to more neighborhoods, so more random nodes have them as a
neighbor. Preferential attachment never has to be imposed. It is induced
automatically the instant a growth rule takes one step out from a random node to
one of its neighbors. The clustering hierarchy has the same local root --- a rule
that makes one link to a node tends to make another among that node's neighbors
--- and the sign of the degree correlations depends on \emph{which} neighbor gets
chosen. Locality is the common cause, and the three mechanisms are three faces of
it.

\section{Summary}

Biological networks grow by duplication: a node is copied along with its
connections, and the copies then diverge by gaining and losing links. Because a
node gains a link whenever one of its neighbors is duplicated, the rule builds in
an effective ``rich get richer'' and hence a power-law degree distribution, while
duplications of self-interacting proteins generate a clustering hierarchy that
link loss steepens. Coupling divergence to duplication adds a density transition
at $q_e=1/2$ and a multifractal degree distribution born of inheritance, and the
model reproduces the disassortative correlations measured in real protein
networks. Together with the two previous chapters, duplication completes a simple
and general principle: a growth rule that attaches to a \emph{neighbor} of a
randomly chosen node attaches preferentially by degree --- and so preferential
attachment, the clustering hierarchy, and degree correlations are all the generic
harvest of locality.

One consequence of the rule is left for the last chapter, because it is not a
statistic of degrees at all. The identical twin a duplication creates is an exact
symmetry, and Chapter~\ref{ch:quantum} turns that into a theorem about how far a
coherent walker can travel on such a network --- then measures eight real interactomes
and finds them confined more tightly than anything modelled here, though whether it is
duplication doing the confining is a question that chapter has to leave open.
